\documentclass[11pt]{article}
\usepackage{amsmath,amssymb,amsthm}
\usepackage{mathtools}
\usepackage{tikz}
\usepackage{enumitem}
\usepackage[margin=1in]{geometry}

\newtheorem{theorem}{Theorem}
\newtheorem*{theorem*}{Theorem}
\newtheorem{lemma}{Lemma}
\newtheorem{definition}{Definition}
\newtheorem*{definition*}{Definition}
\newtheorem{proposition}{Proposition}
\newtheorem{assumption}{Assumption}
\newtheorem{remark}{Remark}

\title{Math 106 -- Notes: Properties of Feynman-Kac Semigroups and application to the filtering problem }
\author{Ethan Levien}
\date{}

\begin{document}
\maketitle


\subsection*{Feynman-kac semigroups}
Previously we were introduced to the Feynman-Kac formula, which gives us a PDE representation of the path expectations 
\begin{equation}
v(x,t)  = {\mathbb E}_{\mu}\!\left[f(X_t)
\exp\!\left(-\int_0^t \lambda(X_s)\,ds\right)|X_0=x
\right]. 
\end{equation}
The idea what that we can connect this expectation to the backward equation to the SDE
\begin{align}
dZ_t &= \lambda(X_t)dt\\
dX_t &= b(X_t)dt + \sigma(X_t)dW_t
\end{align}
because $Z_t = Z_0+\int_0^t\lambda(X_s)\,ds$. This led to the 
\emph{Feynman-Kac} PDE
\begin{equation}\label{eq:fk}
\partial_t v^f(x,t) = {\mathcal L}v^f(x,t) -\lambda(x) v^f(x,t).
\end{equation}
I've introduced the superscript to make explicit that the solution depends on the initial condition $v^f(x,0) = f(x)$. As we see below, it will be important to have this explicit dependence in our notation.
\begin{remark}
In the literature on Feynman-Kac formula, one often considers starting from an initial distribution $\mu$ instead of $X_0 =x$ and write $v_t(f) \coloneqq {\mathbb E}_{\mu}[v^f(x,t)]$. 
\end{remark}
We saw that we can interpret $v^f(x,t)$ as the expectation with respect to an (unnormalized) density function $n(x',x,t)$: 
\begin{equation}
v^f(x,t) = \int n(x',x,t)f(x')dx'
\end{equation}
where $n(x',x,0) = N_0\delta(x'-x)$ with $N_0$ being the initial number of particles. 

Now define 
\begin{equation}
\gamma^f(x,t) = \frac{v^f(x,t)}{v^1(x,t)}
\end{equation}
We can expression this as an expectation 
\begin{equation}
\gamma^f(x,t) = \int p(x',x,t) f(x')dx',\quad p(x',x,t) = \frac{n(x',x,t)}{\int n(x',x,t)dx'}
\end{equation}
Where $p(x',x,t)$ is exactly the density of surviving particles. 

\subsection*{Time-change and long-term behavior}
Now let us consider how we can get a handle on what the solution to the Feynman-Kac PDE looks like. To do so, it is helpful to transform the time axis in a way that makes the death events independent of $X_t$. Let
\begin{equation}
T_t \coloneqq Z_t - Z_0  = \int_0^t\lambda(X_s)ds
\end{equation}
at the new time variable and set $\hat{X}_T = X_t$. Then, since $dt = dT/\lambda(X_t)$, we have
\begin{equation}\label{eq:Xtil}
d\hat{X}_T = \frac{b(\hat{X}_T)}{\lambda(\hat{X}_T)}dT + \frac{\sigma(\hat{X}_T)}{\sqrt{\lambda(\hat{X}_T)}}dW_T
\end{equation}

The generator of this process is ${\mathcal L}/\lambda$ and hence that the Feynman-Kac PDE $\hat{X}_T$ is 
\begin{equation}
\partial_t \hat{v}^f(x,t) = \frac{1}{\lambda(x)} {\mathcal L}v(x,t) -v^f(x,t)
\end{equation}
The operator ${\mathcal L}/\lambda$ and the identity operator commute as do their exponentials $e^{t{\mathcal L}/\lambda}$ and $e^t$. We can therefore express the solution as 
\begin{equation}
\hat{v}^f(x,t) = e^{-t}e^{t{\mathcal L}/\lambda}f(x)
\end{equation}
Since ${\mathcal L}/\lambda$ is the generator of a Markov process, $e^{{\mathcal L}/\lambda} 1 = 1$. It is therefore the case that the long-time behavior of $\gamma^f$, the population expectations is given by the stationary density of the SDE.  




\section*{The filtering problem}

Let $(\Omega,\mathcal F,\mathbb P)$ be a filtered probability space satisfying the usual conditions.  
Let $X_t$ be a time-homogeneous Markov process on $\mathbb R^d$ with generator $\mathcal L$.  
Let $Y_t$ be an observation process taking values in $\mathbb R^d$. We use the usual notation 
\begin{align}
\mathcal F_t^X &= \sigma(X_s:0\le s\le t), \\
\mathcal F_t^Y &= \sigma(Y_s:0\le s\le t), \\
\mathcal F_t &= \sigma(X_s,Y_s:0\le s\le t).
\end{align}
Assume that conditional on the path of $X$, the process $Y$ is Markov with generator $\mathcal H(x)$ acting on functions of $y$.  
Then $(X_t,Y_t)$ is Markov with generator
\begin{equation}
\mathcal A u(x,y)
=
\mathcal L_x u(x,y)
+
\mathcal H(x)_y u(x,y).
\end{equation}
The associated backward equation is
\begin{equation}
\partial_t u(x,y,t)
=
\mathcal L_x u(x,y,t)
+
\mathcal H(x)_y u(x,y,t).
\end{equation}

Given a prior distribution for $X_0$, the filtering problem is to compute
\begin{equation}
\pi_t(f)
=
\mathbb E\!\left[f(X_t)\mid \mathcal F_t^Y\right].
\end{equation}

\subsection*{Poisson observations}

Let $X_t$ be an OU process and let $Y_t$ be given by 
\begin{equation}
Y_t = \Pi\!\left(\int_0^t \lambda(X_s)\,ds\right).
\end{equation}
where $\Pi$ is a unit rate Poisson process. Let $(\Omega, {\mathcal F},\{{\mathcal F}_t\})$ be a probability space which supports both $\{X_t\}$ and $\{Y_t\}$. Let $E$ denote the path space of $X$, hence $X(\cdot): \Omega \to E$ with $X(\omega) = \{X_t(\omega)\}_{0 \le s \le t}$. We let $q$ denote the law of $X$ on $E$. 

For $f\in C^2(\mathbb R)$ the marginal generator of $X_t$ is. 
\begin{equation}
\mathcal L f(x)
=
-\gamma x\, f'(x)
+
\frac{\sigma^2}{2}\, f''(x).
\end{equation}
and conditional generator is
\begin{equation}
\mathcal H(x)\,\varphi(y)
=
\lambda(x)\,\big(\varphi(y+1)-\varphi(y)\big),
\qquad y\in\mathbb N.
\end{equation}
Hence the joint process $(X_s,Y_s)_{s \in [0,t]}$ is generated by 
\begin{equation}
\mathcal A u(x,y)
=
\mathcal L_x u(x,y)
+
\lambda(x)\,\big(u(x,y+1)-u(x,y)\big).
\end{equation}




Let $\{\tau_i\}$ denote the jump times of $Y$, and let
\[
N_t := \max\{ i : \tau_i \le t\}
\]
denote the number of jumps up to time $t$.  
Define the interjump durations
\[
h_i := \tau_i - \tau_{i-1}, \qquad \tau_0:=0.
\]

Conditional on the full path $\{X_s\}$, the survival function of the next interjump time is
\begin{equation}
F_i(h\mid X,\tau_{i-1})
=
\exp\!\left(
-\int_{\tau_{i-1}}^{\tau_{i-1}+h}\lambda(X_z)\,dz
\right).
\end{equation}

For the first $n$ interjump times $(h_1,\dots,h_n)$, the conditional density is 
\begin{align}
\rho_n(h_1,\dots,h_{n}\mid ,{\mathcal F}^X)
&\coloneqq
\prod_{i=1}^{n}
\lambda(X_{\tau_i})
\exp\!\left(
-\int_{\tau_{i-1}}^{\tau_i}\lambda(X_z)\,dz
\right)
\\
&=
\exp\!\left(
-\Lambda_t
+
D_t
\right)\lambda_0^{n}
\end{align}
where $\lambda_0$ is a an arbitrary inverse time-scale to ensure the units makes sense and 
\begin{align}
\Lambda_t  \coloneqq \int_0^{t} \lambda(X_s)\,ds,\quad 
D_t \coloneqq \sum_{i=1}^{N(t)} \log \lambda(X_{\tau_i}) - N(t)\ln \lambda_0 
\end{align}


Let $q$ denote the prior distribution of the full path.  
Conditioning on the observations yields
\begin{equation}
q(X\mid h_1,\dots,h_{N(t)})
=
\frac{
\rho_t(h_1,\dots,h_{N(t)}\mid X)\,q(X)
}{
\mathbb E_q\!\left[\rho_t(h_1,\dots,h_{N(t)}\mid X)\right]
}.
\end{equation}

Thus,
\begin{equation}
\pi_t(f)
=
\frac{
\mathbb E_q\!\left[\rho_t(h_1,\dots,h_{N(t)}\mid X)\,f(X_t)\right]
}{
\mathbb E_q\!\left[\rho_t(h_1,\dots,h_{N(t)}\mid X)\right]
}.
\end{equation}
This provides a means to generate expectations with respect to $X_t$ conditional on the observed process $Y_t$. A major shortcoming of this is that the expectation is highly sensitive to extremal statistics. 

\subsection*{Unnormalized filters and propagators}

Define the unnormalized filtering distribution
\begin{equation}
\gamma_t(f)
=
\mathbb E_q\!\left[\rho_t(h_1,\dots,h_{N(t)}\mid {\mathcal F}_t^X)\,f(X_t)\right],
\end{equation}
so that
\[
\pi_t(f)=\frac{\gamma_t(f)}{\gamma_t(1)}.
\]

For $s<t$, define the unnormalized propagator
\begin{equation}
\Gamma_{s,t} f(x)
\coloneqq
\mathbb E\!\left[
\rho_{s,t}(h_{N(s)+1},\dots,h_{N(t)}\mid X)\,f(X_t) \,\big|\, X_s = x
\right],
\end{equation}
where $\rho_{s,t}$ is the likelihood increment accumulated on $(s,t]$. It follows from the definition that 
\begin{align}
\gamma_t(f)
&= \gamma_s\!\left(\Gamma_{s,t}f\right)\\
\Gamma_{s,t}&= \Gamma_{s,r} \Gamma_{r,t} \qquad (s<r<t).
\end{align}
and therefore 
\begin{equation}
\gamma_t(f)
=
{\mathbb E}[\Gamma_{0,\tau_1}\,
\Gamma_{\tau_1,\tau_2}
\cdots
\Gamma_{\tau_{N(t)},t}f(X_0)],
\end{equation}
with the expectation being over the prior $X_0$.
Now define the scaled interjump propagator
\begin{equation}  
\Gamma_s^{(i)} \coloneqq e^{\lambda_0 s}\Gamma_{\tau_i,\tau_i+ s}. 
\end{equation}
where the factor $e^{\lambda_0}$ will be useful in connecting to population growth. 
Explicitly, this has the form 
\begin{equation}
\Gamma_s^{(i)}f(x) = {\mathbb E}[f(X_{\tau_i + s})e^{\lambda_0 s-(\Lambda_{\tau_i + s}-\Lambda_{\tau_i})}|X_{\tau_i} = x]
\end{equation}
By the Feynmna-Kac fomula, $u_i(x,t) = \Gamma_s^{(i)}f(x)$ is a solution to 
\begin{align}
\partial_t u_i (x,t)&= {\mathcal L}u_i(x,t) + (\lambda_0 - \lambda(x))u_i(x,t)\\
u_i(x,0) &= u_{i-1}(x,h_{i-1})\frac{\lambda(x)}{\lambda_0}
\end{align}
Note that $\lambda_0$ is arbitrary because it will cancel in the ratio. We can therefore select it to be some value such that $\lambda_0 > \lambda(x)$ with high probability. To this end, we assume for simplicity that ${\mathbb P}(\lambda_0> \lambda(X_t))=1$. We will revisit the subtle issues of how to handle deviations from this later on. 
The duel equation 
\begin{align}
\partial_t n_i (x,t)&= {\mathcal L}^*n_i(x,t) + (\lambda_0 - \lambda(x))n_i(x,t)\\
n_i(x,0) &= n_{i-1}(x,h_{i-1})\frac{\lambda(x)}{\lambda_0}
\end{align}
can be viewed as the number density of particles which grow at a rate $\lambda_0 - \lambda(x)$ and are diluted at a $x$ dependent factor $\lambda(x)/\lambda_0$ at the jump times. 


We can then view this from a population perspective as a growth dilution dynamics or a branching OU process. This process should be less sensitive to extremal statistics. 

In particular, consider the branching process consisting of a set $S_t$ particles at time $t$ and for $u \in S_t$ let $X_t^u$ be the position particle $u$. Particles branch at a rate $\lambda_0-  \lambda(x)$. 


\end{document}